% c02.fig.vielecke — the same pair of states in three state spaces in one plane
% (levels C1, C2, C3 of evidence/derivations/c02/golden.json; DERIVATION.md 4.6).
% States: x_1 = (1,0), x_2 = (-3/4, 1/4), priors (1/2, 1/2).
%
% (a) Square with unit circumradius ("diamond"): vertices (1,0), (0,1), (-1,0), (0,-1).
%     x_2 lies on the edge r_x - r_y = -1.  Optimal extreme effect f = (1 + r_x - r_y)/2:
%     f(x_1) = 1, f(x_2) = 0, P = 1.  Level lines r_x - r_y = 2f - 1: f = 0 is the edge
%     through (-1,0) and (0,1), f = 1 the edge through (1,0) and (0,-1).
% (b) Regular pentagon with unit circumradius, vertices (cos(72 k), sin(72 k)) deg.
%     c = cos(36 deg) = (1+sqrt 5)/4 = 0.8090.  Optimal extreme effect
%     f = (c + r_x)/(1 + c): f = 1 on r_x = 1 (through the vertex x_1), f = 0 on
%     r_x = -c, the edge opposite to that vertex.  f(x_2) = (c - 3/4)/(1 + c) > 0,
%     P = (55 - 7 sqrt 5)/40 < 1.  Level lines are vertical.
% (c) Unit disc.  Optimal extreme effect f = (1 + n.r)/2, n = (7,-1)/sqrt(50)
%     = (0.9899,-0.1414).  f = 1 touches the circle at n, f = 0 at -n.
%     P = (8 + 5 sqrt 2)/16 < 1.  Level lines n.r = c for c = -1, 0, 1 and through
%     x_1 (c = 0.9899) is omitted to keep the panel readable; through x_2: c = -0.7778.
% Greyscale: states filled discs; effect level lines dashed; the lines f = 0, f = 1 are
% labelled.
\begin{tikzpicture}[x=1.1cm, y=1.1cm]
  % ---------- (a) ----------
  \begin{scope}
    \draw[tpc axis, fill=tpc-grey-1] (1,0) -- (0,1) -- (-1,0) -- (0,-1) -- cycle;
    \begin{scope}
      \clip (-1.2,-1.2) rectangle (1.2,1.2);
      \foreach \c in {-1,0,1} \draw[tpc effect] ({-1.2},{-1.2-\c}) -- ({1.2},{1.2-\c});
    \end{scope}
    \fill[tpc-state] (1,0) circle (2pt);
    \fill[tpc-state] (-0.75,0.25) circle (2pt);
    \node[tpc label, anchor=west] at (1.05,0.02) {$x_1$};
    \node[tpc label, anchor=south east] at (-0.76,0.27) {$x_2$};
    \node[tpc label, anchor=south] at (-0.55,1.18) {$f = 0$};
    \node[tpc label, anchor=north] at (0.55,-1.18) {$f = 1$};
    \node[tpc label, align=center] at (0,-2.2) {(a) \tpcifenglish{square}{Quadrat}\\[1pt]$P = 1$};
  \end{scope}
  % ---------- (b) ----------
  \begin{scope}[shift={(3.3,0)}]
    \draw[tpc axis, fill=tpc-grey-1] (0:1) -- (72:1) -- (144:1) -- (216:1) -- (288:1) -- cycle;
    \foreach \c in {-0.8090,-0.3568,0.0955,0.5477,1} \draw[tpc effect] (\c,-1.2) -- (\c,1.2);
    % f = 0, 1/4, 1/2, 3/4, 1 at r_x = f (1 + c) - c
    \fill[tpc-state] (1,0) circle (2pt);
    \fill[tpc-state] (-0.75,0.25) circle (2pt);
    \node[tpc label, anchor=west] at (1.05,0.02) {$x_1$};
    \node[tpc label, anchor=south west] at (-0.74,0.27) {$x_2$};
    \node[tpc label, anchor=south] at (-0.809,1.18) {$f = 0$};
    \node[tpc label, anchor=north] at (1,-1.18) {$f = 1$};
    \node[tpc label, align=center] at (0,-2.2)
      {(b) \tpcifenglish{pentagon}{Fünfeck}\\[1pt]$P = \tfrac{55 - 7\sqrt{5}}{40}$};
  \end{scope}
  % ---------- (c) ----------
  \begin{scope}[shift={(6.6,0)}]
    \draw[tpc axis, fill=tpc-grey-1] (0,0) circle (1);
    \begin{scope}
      \clip (-1.2,-1.2) rectangle (1.2,1.2);
      % lines n.r = c: c*n + t*(0.1414, 0.9899)
      \foreach \c in {-1,-0.7778,0,1}
        \draw[tpc effect] ({0.9899*\c - 1.6*0.1414},{-0.1414*\c - 1.6*0.9899})
                       -- ({0.9899*\c + 1.6*0.1414},{-0.1414*\c + 1.6*0.9899});
    \end{scope}
    \fill[tpc-state] (1,0) circle (2pt);
    \fill[tpc-state] (-0.75,0.25) circle (2pt);
    \node[tpc label, anchor=west] at (1.08,0.04) {$x_1$};
    \node[tpc label, anchor=south west] at (-0.72,0.3) {$x_2$};
    \node[tpc label, anchor=south] at (-0.82,1.18) {$f = 0$};
    \node[tpc label, anchor=north] at (0.82,-1.18) {$f = 1$};
    \node[tpc label, align=center] at (0,-2.2)
      {(c) \tpcifenglish{disc}{Scheibe}\\[1pt]$P = \tfrac{8 + 5\sqrt{2}}{16}$};
  \end{scope}
\end{tikzpicture}
